\section{Related Work}
\label{sec:related}

\subsection{Text-to-Image (T2I)}

Denoising diffusion models construct a stochastic \cite{sohldickstein15,songScorematching19,hoDDPM20} or deterministic \cite{DDIM} mapping between the data space and a corresponding-dimension multivariate Gaussian. Signals are synthesized by sampling from a normal distribution and performing a sequence of denoising steps. A number of works \cite{GLIDE,nichol2021improved,DALLE2,Saharia2022PhotorealisticTD} have performed T2I synthesis using images conditioned on the text embedding from a model such as CLIP \cite{CLIP}. Performance is significantly improved in the Latent Diffusion Model \cite{RombachStablediffusion21} (LDM) by performing the diffusion computation in the latent space of a carefully trained variational autoencoder. LDM was trained with a large scale dataset, resulting in the widely adopted Stable Diffusion (\SD) system. We omit the basic diffusion derivation as tutorials are available, e.g., \cite{blogLillog}.

Despite the success of image generation using SD, it is widely acknowledged that SD lacks controllability in synthesis. SD faces challenges in synthesizing multiple objects, often resulting in missing objects or incorrect assignment of prompt attributes to different objects. Recently, ControlNet \cite{controlnet} and T2I-Adapter \cite{mou2023t2i} introduced additional fine-tuning layers to train the model with various forms of image conditioning such as edge maps, or rigging skeletons.

The methods of \cite{layout2im,L2Igan,L2Ixformer,ma2023directed,boxdiff,multidiffusion} have addressed the layout-to-image (L2I) issue using few-shot learning. Directed Diffusion \cite{ma2023directed}, BoxDiff \cite{boxdiff}, and MultiDiffusion \cite{multidiffusion} use coarse bboxes to control subject position, achieving good results by manipulating the spatial latent and text embeddings cross attention map \cite{PromptToPrompt}.

\subsection{Text-to-Video (T2V)}

Text-to-video (T2V) synthesis is generally more difficult than T2I due to the difficulty of ensuring temporal consistency and the requirement for a large paired text and video dataset. \cite{ho2022video,harvey2022flexible,Infilling2022,voleti2022MCVD,yang2022diffusion,nvidianoiseprior} show methods that build on top of image diffusion models. Some works \cite{blattmann2023videoldm,luo2023videofusion} also introduce 3D convolutional layers in the denoising UNet to learn temporal information. Imagen Video \cite{Ho2022ImagenVH} achieves higher resolution by computing temporal and spatial super-resolution on initial low resolution videos. VideoLDM \cite{blattmann2023videoldm} and ModelScope \cite{luo2023videofusion} insert a temporal attention layer by reshaping the latent tensor. Text2Video-Zero \cite{text2video-zero}, \KMtext{denoted as T2V-Zero}, and FateZero \cite{qi2023fatezero} investigate how the temporal coherence can be improved by cross frame attention manipulation with pre-trained T2I models. \cite{nvidianoiseprior} addresses the same problem by introducing temporal correlation in the diffusion noise.
However, these pioneering studies generally lack position control in the video synthesis.

Recently several works have been proposed to solve the controllability in video synthesis problem by using pre-trained models together with low-level conditioning information such as edge or depth maps. Control-A-Video \cite{chen2023controlavideo} and MagicProp \cite{yan2023magicprop} use depth maps with ControlNet to train a temporal-aware network. T2V-Zero \cite{text2video-zero} partially achieves controllability by initializing the latent frames conditioned on the first frame with applied linear translation. However, the control is indirect and requires two steps. First, the user first needs to locate the subject's numerical location, and then adjust a translation offset. Distinct from the methods above, we use an attention injection method to guide the denoising path rather than optimization, and in general this is robust to different random seeds. The recent project Peekaboo \cite{jain2023peekaboo} is concurrent with TrailBlazer and shares similar goals.  Both Peekaboo and TrailBlazer guide subjects in video by manipulating the attention, however the formulations differ in many details. Peekaboo's use of an infinite negative attention injection in the background regions appears to often result in backgrounds with missing detail. 
%Different from our work, Peekaboo uses infinite negative attention injection on the background components. 
In \SecRef{sec:exp}, we will provide both quantitative and visual evidence to demonstrate the better controllability and quality of our results. 
% jp: this is said in the intro: Note that our method not only enables the positioning and animating of bounding boxes through keyframing for subject control, but also allows the manipulation of subject attributes, such as changing the prompt word.
Other very recent preprints that address the layout-to-video problem in differing ways include \cite{lian2023llmgrounded, wang2024boximator, yang2024directavideo}. We do not compare against these concurrent works because their source was not available at the time of writing.
